\documentclass[11pt]{report}

% Packages
\usepackage{amsmath,amssymb,amsthm}
\usepackage{geometry}
\usepackage{cite}
\usepackage{booktabs}
\usepackage{array}
\usepackage{tabularx}
\usepackage{enumitem}
\usepackage{lineno}
\usepackage{xcolor}
\usepackage{hyperref}
\usepackage{tikz}
\usepackage{caption}

% Layout
\geometry{margin=1in}
\newcolumntype{Y}{>{\raggedright\arraybackslash}X}

% Color helpers
\newcommand{\gc}[1]{\textcolor[rgb]{1.00,0.00,0.00}{#1}} % red (gov-comment)
\newcommand{\blue}[1]{\textcolor[rgb]{0.00,0.00,1.00}{#1}} % blue (price)

\title{Spectrum Sensing System Design}
\author{}
\date{}

\begin{document}

\maketitle

\section*{\begin{center} Spectrum Sensing System Design \end{center}}

\begin{description}

\item[Selected SDR:] 

This document focuses on a network-attached SoC option (AD9361 + Zynq-7020 family) and the supporting system design. For comparison only, note that host-attached USB 3.0 SDRs and other form factors exist; they are intentionally not included in the bill-of-materials (BOM) or core hardware sections of this report to avoid vendor-specific BOM entries. The design below therefore centers on the network-attached 7020-class device and generic host compute.

\bigskip

\textit{7020-SDR (AD9361 + Xilinx Zynq-7020):}

\begin{itemize}
    \item Covers $\sim$70\,MHz to 6\,GHz (VHF, UHF, C-band).
    \item Integrated Xilinx Zynq-7020 SoC (Dual ARM + FPGA) for real-time DSP offload.
    \item Gigabit Ethernet interface (network-attached, libiio compatible) — enables long-distance remote deployment and PoE possibility.
    \item Fully compatible with \texttt{libiio}, \texttt{PyADI-IIO}, GNU Radio, and SigMF.
    \item Price examples: \blue{180} \href{https://opensourcesdrlab.com/products/opensourcesdrlab-new-7020-sdr-ad9363-for-pluto-software-defined-radio-without-pa}{\gc{$\bullet$}} -- \blue{220} \href{https://www.ebay.com/}{\gc{$\bullet$}}
\end{itemize}

\bigskip

\textit{Recommended compute node: Jetson Orin Nano / Jetson-class device:}

\begin{itemize}
    \item GPU acceleration (e.g., 128 CUDA cores) for parallel tasks (FFT, ML).
    \item Supports TensorFlow / PyTorch for signal classification and anomaly detection.
    \item Compact and energy-efficient for edge or rooftop gateway deployments.
    \item Price example: \blue{269} (275 TOPS). \href{https://www.ebay.com/itm/168106383528}{\gc{$\bullet$}}
\end{itemize}

\bigskip

\textit{Antenna System Configurations:}
\begin{itemize}
    \item \textit{Channel 1:} VHF/UHF Antenna (70\,MHz -- 3\,GHz)
    \item \textit{Channel 2:} C-Band Antenna (4\,GHz -- 6\,GHz)
\end{itemize}

\begin{table}[h]
    \centering
    \begin{tabular}{@{}lllll@{}}
        \toprule
        \textbf{Band} & \textbf{Antenna Type} & \textbf{Model Example} & \textbf{Frequency Range} & \textbf{Gain} \\ 
        \midrule
        VHF/UHF & Discone & Comet DS-150B & 25\,MHz -- 1.3\,GHz & 2--3\,dBi \\
        VHF/UHF & Log-Periodic & HyperLOG 7060 & 700\,MHz -- 6\,GHz & 5--7\,dBi \\
        C-Band & Horn Antenna & Pasternack PE9851 & 4--8\,GHz & 10--15\,dBi \\
        C-Band & Helical & Custom/DIY & 4--6\,GHz & 12--14\,dBi \\
        \bottomrule
    \end{tabular}
    \caption{Antenna selection for multi-antenna system}
\end{table}

\gc{Recommendation:} add a wideband log-periodic (e.g., HyperLOG 4080, 400\,MHz--8\,GHz). Estimated accessory costs: \blue{476} + SMA 5\,m cable \blue{106} + SMA-to-N adapter \blue{28}. \href{https://www.google.com}{\gc{$\bullet$}}

\newpage

\textit{Supporting Hardware}

\begin{enumerate}[label=\roman*)]
    \item \textbf{Low Noise Amplifiers (LNAs)}
    \begin{table}[!h]
        \centering
        \small
        \begin{tabular}{llcc}
            \toprule
            \textbf{Model} & \textbf{Range} & \textbf{Gain} & \textbf{NF} \\ 
            \midrule
            Nooelec SAWbird+ \blue{85} \href{https://www.amazon.com/Nooelec-SAWbird-ADS-B-Dual-Channel-Applications/dp/B0CBD1KK3G}{\gc{$\bullet$}} & Multi-band (variant) & 20\,dB & $<$1\,dB \\
            Mini-Circuits ZX60-33LN-S+ \blue{159} \href{https://www.ebay.com/itm/136591156802}{\gc{$\bullet$}} & 50\,MHz -- 3\,GHz & 20\,dB & 3\,dB \\
            SparkGap LNA4ALL & 28\,MHz -- 2.5\,GHz & 20\,dB & 1.5\,dB \\
            \bottomrule
        \end{tabular}
    \end{table}
    \smallskip
    Note: SAWbird+ variants include band-specific SAW filters; select the correct variant per band or use wideband LNAs with external filters.
    
    \item \textbf{Automatic Antenna Switching System}
    \begin{itemize}
        \item RF switch examples: Mini-Circuits ZSWA-4-30DR+ (DC--3\,GHz, 4-port) or Koarmy QM-SW-4S (0.1--6\,GHz). Control via Jetson GPIO or a small MCU.
        \item Preselection (protect front-end / prevent overload): FM broadcast notch (88--108\,MHz), bandpass tiles for 700--960\,MHz, 1.7--2.7\,GHz, 3--4\,GHz, 4--6\,GHz.
        \item Stable timing/clocking: GPSDO or external 10\,MHz + PPS for frequency accuracy and multi-node coherence (TDOA).
        \item Digitally-controlled step attenuator (0--31\,dB) to avoid ADC saturation in strong-signal environments.
    \end{itemize}
    
    \item \textbf{Outdoor / Rooftop Monitoring}
    \begin{itemize}
        \item Use non-penetrating roof mounts or standard mast mounting (1.25''--2'').
        \item Weatherproof enclosures and grounding / lightning protection.
        \item Low-loss coax (LMR-400) for longer cable runs to minimize insertion loss.
        \item For network-attached SDRs, use shielded Cat6 and PoE to simplify rooftop deployments (100\,m runs without active USB repeaters).
    \end{itemize}
\end{enumerate}

\gc{Practical recommendation:}
\begin{itemize}
    \item LNAs: Nooelec SAWbird+ for narrow-band needs or Mini-Circuits ZX60 for wideband low-NF amplification.
    \item Interface: prefer network-attached SDRs (GigE / libiio) for remote rooftop installations; host-attached USB 3.0 SDRs are viable for local/short-run deployments but are excluded from BOM entries here.
\end{itemize}

\textit{Testing \& Calibration}

\begin{description}
    \item[libiio-utils:] Use \texttt{iio\_info} and \texttt{iio\_readdev} to validate device reachability and streaming for network-attached SDRs.
    \item[VSWR / Antenna:] Use a NanoVNA or equivalent to verify antenna matching.
    \item[Frequency Sweep:] Sweep available frequency ranges to characterize front-end response and identify spurs.
    \item[Noise Floor:] Measure baseline noise levels for sensitivity and detection threshold tuning.
\end{description}

\bigskip

\textit{Software Tools}

\begin{description}
    \item[GNU Radio:] Flowgraph development, spectrum estimation, demodulation blocks.
    \item[GQRX:] Quick spectrum visualization and tuning.
    \item[URH (Universal Radio Hacker):] Protocol analysis and decoding.
    \item[PyADI-IIO / libiio:] For Python control and streaming with AD9361-based network SDRs.
    \item[Vendor / host drivers:] Use appropriate host drivers or vendor-provided libraries for device control and streaming.
\end{description}

\bigskip

\emph{Compute and processing pipeline}

\begin{enumerate}[leftmargin=*,label={\arabic*.}]
    \item \gc{ML pipeline:} Use PyTorch/TensorFlow on Jetson (or Orin) for classification; leverage cuFFT/cuSignal where available.
    \item \gc{Data transport:} Prefer 8-bit IQ for streaming/storage to reduce bandwidth unless dynamic range requirements force higher resolution captures.
    \item For host interfaces, optimize host-side buffers and data paths (whether Ethernet or USB) and choose moderate sustained sample rates (e.g., 2--10\,Msps per channel) for reliable real-time operation.
    \item For network-attached 7020-class SDRs, offload decimation/filtering to the device SoC/FPGA and tune libiio buffer settings to prevent dropped samples.
    \item \gc{UI / orchestration:} Jetson host runs \texttt{iiod} client(s) and a Flask (or Django) web dashboard with gr-specest waterfalls and REST APIs.
\end{enumerate}

\bigskip

\emph{Upgrade cost matrix and example budget}

\begin{table}[h]
    \centering
    \begin{tabular}{@{}l p{7.2cm} r@{}}
        \toprule
        \textbf{Item} & \textbf{Benefit} & \textbf{Cost (\$)} \\
        \midrule
        7020-SDR (AD9361) & Standalone SoC architecture, GigE & 200 \\
        Jetson Orin / Nano & ML / DL acceleration & 269 \\
        Sync cable / GPSDO & Coherent timing and PPS/10 MHz refs & 200 \\
        Antenna system (Comet + HyperLOG 7060) & Wideband reception & 1,075 \\
        RF networking & Switch, LMR-400, Cat6, PoE & 250 \\
        LNAs \& filters & Noise floor optimization & 150 \\
        \textbf{Estimated total (example)} & & $\sim$ \blue{\$1,944--3,000} \\
        \midrule
        RFSoC node & High-bandwidth FPGA/DAC (optional) & \gc{+5,000} \\
        \bottomrule
    \end{tabular}
    \caption{Example upgrade / procurement cost matrix}
\end{table}

\vspace{6pt}

\item[Similar projects \& references:] Representative examples and resources:
\begin{enumerate}[label=\roman*)]
    \item \emph{DARPA Spectrum Collaboration Challenge (SC2)}: cognitive radio / spectrum-sharing research.
    \item \emph{University testbeds}: Berkeley, Virginia Tech and others provide SDR testbeds and cognitive-radio research.
    \item \emph{Commercial solutions}: ThinkRF, Per Vices (spectrum monitoring appliances).
    \item \emph{Open-source projects:} GNU Radio, OpenAirInterface (OAI), SigMF (Signal Metadata Format), Pothos SDR, gr-specest, and various GitHub repositories with spectrum scanning and ML classifiers.
\end{enumerate}

\vspace{8pt}

\noindent\textbf{Throughput note (comparison):} Practical host-interface throughput depends on the chosen hardware: typical host-limited USB 3.0 practical throughput is on the order of a few hundred MB/s; for very high-throughput needs prefer 1/10\,GbE-capable platforms and appropriate network interfaces. This is a general comparison note only — host-attached USB SDRs are intentionally not itemized in the BOM/hardware lists above.

\vspace{8pt}

\noindent\textbf{Final remarks (content parity):}  
All material has been normalized around the network-attached 7020-class SDR and the Jetson compute node. Mentions of host-attached USB SDRs have been removed from BOM and other hardware sections; a short comparison note is included to preserve guidance for evaluation without creating vendor-specific BOM entries.

\end{description}

\end{document}
